\documentclass[letterpaper,11pt,leqno]{article}
\usepackage{paper}
\bibliographystyle{paper}

% Enter paper title to populate PDF metadata:
\hypersetup{pdftitle={Velocity Verlet in 3D}}

% Enter path to BibTeX file with references:
\newcommand{\bib}{paper.bib}

% Enter path to PDF file with figures:
\newcommand{\pdf}{figures.pdf}

\begin{document}

\begin{titlepage}

% Enter title:
\title{Velocity Verlet in 3-D}
% Enter authors:
\author{Guy Matz}
% Enter affiliations and acknowledgements:
% \thanks{Guy Matz: Chemistry Department, University of Padua}

% Enter date:
% \date{2026-07-09}
% Enter paper URL (can be commented out):
% \paperurl{https://github.com/guymatz/guymatz-cc2526/verlet/verlet3/paper}
\maketitle
% Abstract:
The ``Velocity Verlet'' algorithm is implemented in order to predict the location of three
particles in a system
\\~\\
\\~\\
Summarize all aspects of the report, including purpose, method, results, conclusions, and recommendations. There are no graphs, charts, tables, or images in an abstract. Similarly, an abstract does not include a bibliography or references.
\\~\\
Highlight important discoveries or anomalies. It's okay if the experiment did not go as planned and
necessary to state the outcome in the abstract.
\end{titlepage}

%% Introduction
\section{Introduction}\label{s:introduction}
In the field of Chemistry, predicting how individual molecules will move is crucial to 
understanding the interaction of molecules in a system.  This gives the chemist a view
of the evolution of a system and can be helpful in predicting dynamical and structural 
properties of a collection of interacting molecules.
\\~\\
Newton's equations of motion have historically been used to estimate the trajectories of 
interacting molecules.  One might use a method such as Euler's for determining numerical 
solutions to Newton's equations, however in the last century a new method was 
discovered by the French physicist Loup Verlet.  The so-called ``Verlet Integration'' provides
a more accurate method for numerical integration of equations of motion.
\\~\\
An implementation of the Verlet algorithm in Fortran is included in the supplemental 
materials.  With it we are able to determine the location of three molecules in 3D space.
\\~\\
The JPCA15 algorithm\footnote{Rampino S, The Journal of Physical Chemistry A 120, 4683--4692,
2016} gives us the mechanism for computing the ``potential energy and the vector of the
derivatives of the potential with respect to the inter---atomic distances''.
The Lennard-Jones algorithm is typically used to calculate potential energy, however here
we use the JPCA15 algorithm as it allows us to compute potentials for three interacting
particles, rather than the two allowed for by Lennard-Jones.
We use these results as input to our Velocity Verlet algorithm.

One thing to note: Numerical integrations are, by nature, inaccurate.  Errors accumulate \dots
 
The
literature\footnote{https://pythoninchemistry.org/ch40208/molecular\_dynamics/choosing\_the\_right\_timestep.html}
recommends using a timestep of ``one-tenth of the period of the fastest vibration in your
system.  For systems containing hydrogen atoms, this typically means using timesteps of around
0.5 to 1.0 femtoseconds''.  Our system contains only hydrogen atoms, so we proceed with a 
timestep of 0.5 femtoseconds.

\section{Experimental Section}\label{s:experiment}

The initial setup of the system is three hydrogen atoms --- $A, B, \&\ C$ --- in 3D space,
arranged along the x-axis as follows:
\begin{itemize}
  \item $H_A$ @ (-10, 0, 0)
  \item $H_B$ @ (10, 0, 0)
  \item $H_C$ @ (10.7, 0, 0)
\end{itemize}
We experiment with different `initial velocities' in order to find at least one non-reactive
and one reactive trajectory.  We are working solely in the $x$ ``direction'' so we can ``play''
with the velocity of atom $A$ in our setup.  We start with an initial velocity of 0.1 for
$A_x$, and increase to 1.0 by steps of 0.1.

\subsection{Subsection with references} 

% References can be appended at the end of a sentence, in parenthesis \citep{MS15}. References can also include page numbers or other bibliographical information \citep[p. 1305]{MS19}.  References can be in text: for instance, \citet{M12} found this and \citet[figure 1]{M14} found that. It's also possible to collate several references to the same author: for instance, \citet{M12,M14} found things. References from the same authors and same year appear as follows: \citet{LMS18a,LMS18b} wrote something. It is also possible to cite the authors of a study by name, or the year of a study: for instance, \citeauthor{EMM21} wrote a paper in \citeyear{EMM21} on this topic.


\section{Results}\label{s:results}

This section displays a number of mathematical expressions to showcase the math fonts used in the template.

\subsection{Some equations} 

\section{Discussion}\label{s:discussion}

Here is a section with a variety of graphs. At vero eos et accusamus et iusto odio dignissimos ducimus, qui blanditiis praesentium voluptatum deleniti atque corrupti, quos dolores et quas molestias excepturi sint, obcaecati cupiditate non provident, similique sunt in culpa, qui officia deserunt mollitia animi, id est laborum et dolorum fuga. 

\subsection{Subsection with graphs at the top of the page}

A simple two-panel graph is on figure \ref{f:graph1}. It will be placed at the top of the page, just about here. Et harum quidem rerum facilis est et expedita distinctio. Nam libero tempore, cum soluta nobis est eligendi optio, cumque nihil impedit, quo minus id, quod maxime placeat, facere possimus, omnis voluptas assumenda est, omnis dolor repellendus. Temporibus autem quibusdam et aut officiis debitis aut rerum necessitatibus saepe eveniet, ut et voluptates repudiandae sint et molestiae non recusandae.

\begin{figure}[t]
\subcaptionbox{A first panel, 1951--2019\label{f:panelA}}{\includegraphics[scale=0.2,page=1]{\pdf}}\hfill
\subcaptionbox{A second panel, 1951--2019\label{f:panelB}}{\includegraphics[scale=0.2,page=2]{\pdf}}
\caption{Graph with two panels}
\note{This is a note for the graph. Nam libero tempore, cum soluta nobis est eligendi optio, cumque nihil impedit, quo minus id, quod maxime placeat, facere possimus.}
\label{f:graph1}\end{figure}

\subsection{References to figures and panels} 

\section{Conclusion}\label{s:conclusion}

% As usual with LaTeX, it is easy to refer to a figure: see figure \ref{f:graph1}. It is possible to refer to a specific panel in a figure, for instance figure \ref{f:panelA} or figure \ref{f:panelB}. Its also possible to refer to the entire figure, for instance figure \ref{f:graph1} or figure \ref{f:graph2}.

\end{document}
